\documentclass[8pt]{strategynet-journal}
\usepackage{volt-paper-style}
\hypersetup{pdftitle={Options, Exercise and Learned Hedging: A Practical Primer},pdfauthor={research@strategynet.ai}}
\newcommand{\E}{\mathbb{E}}
\newcommand{\Q}{\mathbb{Q}}
\newcommand{\PP}{\mathbb{P}}
\newcommand{\pos}[1]{\left(#1\right)^{+}}
\SNHeader{strategynet.ai:Options and learned hedging}
\title{Options, Exercise and Learned Hedging:\\A Practical Primer}
\author{\href{mailto:research@strategynet.ai}{research@strategynet.ai}}
\date{October 2026}
\begin{document}
\maketitle
\begin{abstract}
\normalsize An institution that sells an option receives money today and takes on a payment whose amount, and sometimes timing, depend on future markets. Stock positions can offset some of that exposure, while rebalancing incurs costs and early exercise can bring settlement forward. Understanding a learned hedge therefore begins with the contract, its financed trading account and the market conditions under which decisions are assessed. This primer develops the background needed to understand the option-pricing and hedge-learning applications in Volt and Jormungandr. We begin with a put option, follow its cash flows, introduce simulated market paths and distinguish forecasting from valuation. We then derive the holder's exercise decision and the writer's hedging problem, explain how training targets and derivatives are obtained, and connect the mathematics to the application interfaces. Worked calculations, compact diagrams and review questions accompany the discussion. Completed one-stock experiments provide a practical reference, while a reading guide leads into the original Deep Hedging and learned optimal stopping papers. The assumed background is basic algebra, functions and the idea of a derivative; probability, stochastic models and learning terminology are introduced as they become useful.
\end{abstract}
\SNtwocolumntableofcontents
\begin{SNcolumns}

\section{Options and contract terms}
Suppose a fund owns shares and wants protection against a fall in their value.
It buys a put: the right to sell a specified amount of stock at an agreed
\emph{strike price}, subject to the contract's exercise terms. The seller,
often called the writer, receives a \emph{premium} for accepting the obligation.
The fund now has protection, and the writer has a risk to manage.

Write $S_t$ for the stock price at time $t$, $K$ for the strike and $T$ for
expiry. A cash-settled put pays $g(S_T)=(K-S_T)^+$ at expiry, where
$x^+=\max(x,0)$. A call pays $(S_T-K)^+$. With $K=100$, a put pays 10 if
the stock finishes at 90 and zero if it finishes at 110. These payments are
per underlying unit. The contract multiplier and position size turn them
into the portfolio's actual currency payments.

An option's \emph{payoff} describes the payment when it settles. Its
\emph{value} before settlement is the amount assigned to its remaining rights
by a market quote or a valuation model. A put with the stock at 100 may be
valuable because the stock can fall before expiry. Its immediate exercise
payment is currently zero. Time, uncertainty, dividends and financing all
enter the valuation~\cite{oicpricing}.

\begin{center}\begin{minipage}{\linewidth}\centering
\begin{tikzpicture}[x=.070cm,y=.065cm,font=\normalsize]
\draw[->] (55,0)--(148,0) node[right]{$S_T$};
\draw[->] (60,0)--(60,49) node[above]{Payment};
\draw[darkblue,thick] (60,40)--(100,0)--(143,0);
\draw[darkgreen,thick,dashed] (65,0)--(100,0)--(143,43);
\node[darkblue,anchor=west] at(64,37){Put};
\node[darkgreen,anchor=west] at(119,31){Call};
\draw (100,0)--(100,-2) node[below]{$K=100$};
\draw (60,40)--(58,40) node[left]{40};
\end{tikzpicture}
\captionof{figure}{Terminal payments for a call and put with strike 100.
The premium is paid separately when the option is purchased.}
\end{minipage}\end{center}

Exercise rights determine when the holder can take the payment. A
\emph{European} option permits exercise at expiry. A \emph{Bermudan} option
specifies a finite set of dates. An \emph{American} option permits exercise
through its contractual interval. In numerical work, increasingly frequent
exercise dates can approximate continuous exercise. Physical settlement
also specifies delivery of stock and strike cash; our worked learning
examples use the economically specified cash payment.

The holder and writer have separate decisions. The holder chooses when to
exercise. The writer chooses what to hold while the obligation is alive.
A stock purchase can offset a short call's exposure to a rise in price; a
stock short can offset a short put's exposure to a fall. Rebalancing costs
money, so the useful hedge balances protection and execution costs. The
exercise rule determines when that trading process ends.

We will follow a one-share put throughout. Illustrative calculations use
explicitly chosen prices and payments. The measured training results appear
in Section~\ref{sec:evidence}, where the model, units and sample sizes are stated.

\section{Financing, units and settlement}
A hedge is a financed portfolio. If a writer receives premium $p$ and chooses
$h_0$ stock units at price $S_0$, its cash balance becomes
\[
 b_0=p-h_0S_0-\lambda S_0|h_0|,
\]
where $\lambda$ is a proportional stock trading fee. A negative $h_0$ is a
stock short and raises cash when opened. Negative cash represents borrowing.
At the next date, cash earns or pays interest and the stock is marked at its
new price. The combined value is the money available to meet the liability.

Let $b_n$ be cash after trading at date $t_n$, and let $D_{n+1}$ be a cash
dividend per share paid at the next date. With constant annual funding rate
$r$ and interval $\Delta t_n=t_{n+1}-t_n$, the pre-trade wealth on arrival, for a continuing contract with $S_{n+1}$
measured ex-dividend, is
\[
 w_{n+1}^{-}=b_ne^{r\Delta t_n}+h_n(S_{n+1}+D_{n+1}).
\]
Choosing new holdings $h_{n+1}$ then leaves
\[
 b_{n+1}=w_{n+1}^{-}-h_{n+1}S_{n+1}
             -\lambda S_{n+1}|h_{n+1}-h_n|.
\]
This formula pays for the inventory adjustment and its fee. It describes a
self-financing hedge: after the initial premium, trades use the existing
stock and cash account. More elaborate funding arrangements can supply
different rates or explicit borrowing constraints.

For a concrete one-period calculation, sell a put for 6 with stock at 100,
short 0.4 shares, and charge a 0.1\% stock fee. Opening cash is
$6+40-0.04=45.96$. Suppose the stock subsequently reaches 90, cash grows by
1.001, and the put pays 10. Closing the stock short costs 36 and its closing
fee is 0.036. The remaining cash is
\[
 45.96\times1.001-36-0.036-10=-0.03004.
\]
Without the stock trade, the corresponding result is $6\times1.001-10=-3.994$.
The hedge offsets most of this particular loss. Another stock outcome produces
a different result; the policy will be assessed across many paths.

If exercise occurs at time $\tau<T$, the account closes its hedge there,
pays the contract and keeps the remaining cash invested until $T$. Reporting
all path profits in currency at this common final date makes early and late
settlements comparable. At constant rates, discounted profit is the final
profit multiplied by $e^{-rT}$. A study must say which convention its loss uses.

\begin{center}\begin{minipage}{\linewidth}\centering
\begin{tikzpicture}[font=\normalsize,x=1cm,y=1cm]
\draw[->,thick] (0,0)--(6.8,0);
\foreach \x/\lab in {.35/{$0$},2.5/{$t_1$},4.5/{$\tau$},6.4/{$T$}}
 {\draw (\x,.08)--(\x,-.08) node[below]{\lab};}
\node[align=center] at(.7,.62){Receive premium\\Open hedge};
\node[align=center] at(2.6,1.25){Rebalance\\Pay trading fee};
\node[align=center] at(4.6,.62){Close hedge\\Pay put};
\draw[darkblue,thick] (4.5,-.65)--(6.4,-.65);
\node[align=center,anchor=north] at(5.45,-.72){Cash financing};
\end{tikzpicture}
\captionof{figure}{An early-settled path. Trading ends at exercise; the
reported account continues in cash to the common horizon.}
\end{minipage}\end{center}

Contract units deserve the same attention as the model. If a put is quoted
at 6 per share and its multiplier is 100, one contract costs 600. A model
delta of $-0.4$ means $-40$ shares for that contract's option exposure. For a
writer hedging the short put, holding those $-40$ shares offsets its positive
stock-price sensitivity. Fees follow the actual shares traded. Applying the
multiplier at both the payoff producer and the learner would scale the
obligation twice.

Settlement timing also matters around dividends. If the exercise opportunity
comes just before a cash dividend, the holder observes the pre-dividend stock
price. A continuing stock position receives the dividend and moves to the
ex-dividend mark. A hedge closed before that event receives the appropriate
pre-dividend proceeds. Volt's dated scenario producer records this ordering
with its settlement and income arrays.

\section{Probability and valuation}
The future stock price is uncertain. A \emph{scenario} specifies one possible
sequence of prices, dividends and other observations. A \emph{random variable}
is a number determined by that uncertain future, such as the put's payment.
An expectation $\E[X]$ averages a random variable $X$ using its probability law.
For two equally likely payments of zero and 10, the expectation is 5.

Expectations conditional on current observations are especially useful.
$\E[X\mid S_t=s]$ asks for the average future outcome among situations with
that current stock state. A continuation-value model estimates this average
payment from waiting. Its output changes with the information available now.
A decision at date $t$ uses observations up to that date, which we denote by
$\mathcal F_t$. This growing information set is called a \emph{filtration}.
It is the mathematical expression of observing the market as time passes.

Dispersion describes what an average leaves out. Variance is
$\operatorname{Var}(X)=\E[(X-\E[X])^2]$, and its square root is the standard
deviation. Covariance measures how two variables move together; correlation
rescales covariance by their standard deviations. A return generator for
several stocks must describe these joint movements. Drawing every stock
independently would impose zero dependence whether or not the intended
model has that property.

We use two probability laws for different purposes. The \emph{physical law},
$\PP$, models actual return outcomes and is used for forecasting and
real-world risk assessment. A \emph{risk-neutral law}, $\Q$, is a pricing
construction consistent with the selected tradable assets and financing
assumptions. In the frictionless models used here, discounted gains including
dividends are martingales under $\Q$: their expected next value, conditional
on current information, equals their current value.

For a European payment $Z$ and deterministic rates, its model value is
\[
 V_t=\E^{\Q}[e^{-r(T-t)}Z\mid\mathcal F_t].
\]
The exponential discounts future currency back to date $t$. A stock's physical
expected return can include a risk premium or predictive signal. Under the
GBM pricing specification, its ex-dividend drift becomes $r-q$, where $q$
is the continuous dividend yield. That change in probability law supplies
the appropriate pricing expectation; the forecast retains its separate role.

This distinction becomes practical when we add factors. A factor model might
forecast a stock rise under $\PP$. The writer can then benefit from retaining
directional exposure, depending on its risk preference. A pricing calculation
still needs coherent $\Q$ dynamics or observed option marks. In a market with
transaction costs and incomplete hedges, preferences and execution conditions
also influence acceptable prices and positions. Our controlled pricing tests
state a particular frictionless pricing model before evaluating trading costs.

A \emph{Monte Carlo} calculation replaces an expectation by an average over
simulated paths. For independent path outcomes $X_1,\ldots,X_N$, define
\[
 \bar X=\frac1N\sum_{i=1}^N X_i,\qquad
 \operatorname{SE}(\bar X)=\frac{s_X}{\sqrt N},
\]
where $s_X$ is their sample standard deviation. Under the usual finite-variance
central-limit approximation, $\bar X\pm1.96\operatorname{SE}$ gives a 95\%
interval for the expectation. Four times as many independent paths roughly
halves the standard error. Simulation bias, parameter uncertainty and model
error require additional checks beyond that sampling interval.

When comparing two hedges, evaluate both on the same paths. The difference
on each path then provides a paired observation, often giving a more precise
comparison than two unrelated samples. Training paths improve the model;
validation paths select its settings or checkpoint; fresh assessment paths
measure the selected policy. Keeping these roles separate makes the final
comparison useful. Repeatedly choosing policies using the assessment outcomes
would turn that cohort into another selection dataset.

\section{Market models, prices and Greeks}
A stock model specifies how scenarios evolve. The simplest model used in
our examples is geometric Brownian motion, or GBM:
\[
 dS_t=\mu S_t\,dt+\sigma S_t\,dW_t.
\]
The parameter $\mu$ is proportional drift, $\sigma$ is annualized volatility
when time is measured in years, and $W_t$
is Brownian motion. A Brownian increment over an interval of length $\Delta t$
is Gaussian with mean zero and variance $\Delta t$, independently of earlier
increments. The symbol $dW_t$ represents that stochastic change in the
continuous-time equation. The model assumes continuous positive stock paths
and constant volatility.

GBM has an exact transition over a time interval:
\[
 S_{t+\Delta t}=S_t\exp\!\left[
   (\mu-\tfrac12\sigma^2)\Delta t+\sigma\sqrt{\Delta t}\,Z\right],
 \quad Z\sim N(0,1).
\]
Here $N(0,1)$ denotes a standard normal draw. The $-\sigma^2/2$ term follows
from the stochastic change of variables for $\log S$. Under the pricing law,
$\mu=r-q$. With several stocks, correlated Gaussian draws specify their joint
Brownian increments. Exact GBM transitions remove time-stepping error for the
stock at the simulated dates. The exercise and hedge calendars still determine
which decisions are available.

Under GBM with constant volatility, rates and yield, the European call has
an analytic value. For $\sigma>0$ and remaining maturity $\tau=T-t>0$,
\[
\begin{aligned}
 C&=Se^{-q\tau}\Phi(d_1)-Ke^{-r\tau}\Phi(d_2),\\
 d_1&=\frac{\log(S/K)+(r-q+\sigma^2/2)\tau}{\sigma\sqrt\tau},\qquad
 d_2=d_1-\sigma\sqrt\tau.
\end{aligned}
\]
$\Phi(z)$ is the probability that a standard normal variable is at most $z$.
The formula assumes frictionless trading and the declared continuous GBM model.
The corresponding European put is $P=C-Se^{-q\tau}+Ke^{-r\tau}$, the put--call
parity relation. It follows by comparing portfolios with the same terminal
payments. Call delta is $e^{-q\tau}\Phi(d_1)$; put delta subtracts $e^{-q\tau}$.
These analytic quantities supply useful numerical controls. The early-exercise
problem adds the right to settle before the formula's terminal date.

Heston gives volatility its own changing state. Write $v_t$ for instantaneous
variance. A common pricing specification is
\[
\begin{aligned}
 dS_t&=(r-q)S_t\,dt+\sqrt{v_t}S_t\,dW_t^S,\\
 dv_t&=\kappa(\theta-v_t)\,dt+\xi\sqrt{v_t}\,dW_t^v,\\
 d\langle W^S,W^v\rangle_t&=\rho\,dt.
\end{aligned}
\]
For scale, volatility 20\% corresponds to variance $v=0.2^2=0.04$.
The positive parameter $\kappa$ controls mean reversion toward variance level
$\theta$; $\xi$ controls variance fluctuations; $\rho$ couples stock and variance
shocks. This adds uncertainty that a stock-only hedge can leave exposed.
Numerical simulation must also handle variance near zero and integration of
variance between dates. Our European Heston studies and constant-volatility
early-exercise studies use these distinct model specifications~\cite{heston}.

Once a model assigns option value $V(S,\sigma,r,t)$, its derivatives describe
local changes. Delta is $\Delta=\partial V/\partial S$, gamma is
$\Gamma=\partial^2V/\partial S^2$, vega is $\nu=\partial V/\partial\sigma$,
rho is $\partial V/\partial r$, and theta here is the derivative with respect
to advancing calendar time. A short-horizon approximation is
\[
 \delta V\approx\Delta\,\delta S+
 \tfrac12\Gamma(\delta S)^2+\nu\,\delta\sigma+
 \frac{\partial V}{\partial r}\delta r+\Theta\,\delta t.
\]
It is a local expansion with omitted cross terms and higher-order effects.
A sufficiently large move calls for full repricing of the scenario.

The units determine how to read these numbers. Vega differentiated with
respect to volatility as a decimal must be multiplied by 0.01 for a
one-percentage-point volatility change. Annualized theta needs the declared
day-count conversion for a daily display. Signed quantities and multipliers
aggregate contract Greeks into portfolio Greeks. A delta near zero can coexist
with substantial gamma or vega, so the portfolio's response can still change
rapidly as the market moves.

\emph{Implied volatility} reverses the pricing calculation: choose volatility
so the model price matches an observed option quote. Different strikes and
expiries generally imply different values, producing the volatility grid
commonly called a volatility surface. The fitted volatility is conditional
on the pricing method, exercise rights, rates and dividends. Calibration
errors or stale inputs can resemble trading opportunities in a cross-section.
A practitioner therefore examines quote times and bid--ask prices alongside
the calculated exposures. Quotes and open interest describe market observations
and outstanding contracts; portfolio risk additionally requires signed positions.
A sensitivity to initial Heston variance, $\partial V/\partial v_0$, also
carries different units from a derivative with respect to a scalar volatility.
Each displayed Greek should identify the parameter being perturbed.

\section{Exercise and dynamic programming}
A holder of a put may exercise because the immediate payment is more useful
than retaining the option. On a finite calendar, this becomes a sequence of
comparisons. Let $V_n(s)$ be the option value at date $t_n$ when stock is $s$.
At expiry, $V_N(s)=g(s)$. At an earlier exercise opportunity,
\[
\begin{aligned}
 C_n(s)&=\E^{\Q}[e^{-r\Delta t_n}V_{n+1}(S_{n+1})\mid S_n=s],\\
 V_n(s)&=\max\{g(s),C_n(s)\}.
\end{aligned}
\]
$C_n$ is the \emph{continuation value}. On dates where exercise is unavailable,
the value is the continuation value alone. Starting from the known terminal
payoff and moving backward is \emph{dynamic programming}. Each earlier value
uses the already computed later decisions.

A small tree makes the comparison concrete. Take a put with strike 100 and
stock at 100. Exercise is allowed after either of two intervals. At each
interval the stock multiplies by 1.2 or 0.8, while cash grows by a factor of
1.05. The pricing probability of an up move is chosen so the discounted stock
has unchanged conditional expectation:
\[
 p_{\rm up}=\frac{1.05-0.8}{1.2-0.8}=\frac58.
\]
At the low intermediate stock price of 80, exercise pays 20. Continuing gives
\[
 C_1(80)=\frac{(5/8)4+(3/8)36}{1.05}=15.2381.
\]
The holder exercises. At stock 120, immediate payment is zero, whereas
continuation is $[(5/8)0+(3/8)4]/1.05=1.4286$, so the holder waits. The initial
Bermudan value is $[(5/8)1.4286+(3/8)20]/1.05=7.9932$. The terminal-only
European value on the same tree is 6.2925. The additional exercise opportunity
therefore contributes about 1.7007. The factor 1.05 is the chosen funding
factor per interval; it is distinct from an annual continuously compounded rate.

\begin{center}\begin{minipage}{\linewidth}\centering
\begin{tikzpicture}[font=\normalsize,x=1cm,y=1cm,
 every node/.style={align=center},->,>=Latex]
\node (a) at(0,0){$S=100$\\$V=7.9932$};
\node (b) at(2.55,1.05){$S=120$\\Wait: $C=1.4286$};
\node (c) at(2.55,-1.05){$S=80$\\Exercise: $g=20$};
\node (d) at(5.7,1.75){$S=144$; $g=0$};
\node (e) at(5.7,0){$S=96$; $g=4$};
\node (f) at(5.7,-1.75){$S=64$; $g=36$};
\draw (a)--node[above left]{$5/8$}(b);
\draw (a)--node[below left]{$3/8$}(c);
\draw (b)--(d);\draw (b)--(e);\draw[dashed] (c)--(e);\draw[dashed] (c)--(f);
\node[anchor=south] at(2.8,2.15){Read prices forward; compute values backward};
\end{tikzpicture}
\captionof{figure}{A two-interval teaching tree. Up moves have pricing
probability $5/8$ and down moves $3/8$. Dashed branches show stock outcomes
after the put has already settled at the low intermediate node.}
\end{minipage}\end{center}

Trees and finite-difference grids approximate this recursion directly.
The finite-difference method solves a pricing equation over stock and time;
American exercise adds the condition that continuing value and immediate
exercise satisfy an obstacle problem. Refining the spatial and temporal
meshes checks numerical stability. As the number of underlying assets grows,
a full grid requires many more points: $m$ stock levels for each of $d$ assets
produce $m^d$ combinations. This growth motivates simulation and learned
continuation functions.

A \emph{stopping time} $\tau$ is a rule that decides to stop using the
information available when it stops. For example, ``exercise when immediate
payment exceeds today's estimated continuation'' can be applied as the path
unfolds. Selecting the most profitable exercise date after seeing the whole
future would give the holder extra information. The numerical tests preserve
this chronology both in exercise decisions and policy observations.

\section{Neural policies and training objectives}
A neural network is a parameterized function. Its weights, collected in
$\theta$, transform inputs through layers into an output. A hidden layer
computes an affine transformation $Ax+b$, then applies an activation such
as $\tanh$ or the rectified linear function $\max(x,0)$. Changing the weights
changes the function. Training chooses weights that improve a stated objective
on the supplied examples.

For continuation regression, the input is today's state. A future stopping
rule, already selected during the backward calculation, supplies a discounted
payment on each training path. Use the notation
$G_n=e^{-rt_n}g(S_n)$ and let $\tau_{n+1}$ be that later rule. The regression
fits a function $c_\theta$ by minimizing
\[
 \frac1B\sum_{i=1}^{B}
 \left[c_\theta(S_n^{(i)})-G_{\tau_{n+1}^{(i)}}^{(i)}\right]^2.
\]
The batch contains $B$ paths. Here $c_\theta$ estimates value discounted to
time zero. Multiplying it by $e^{rt_n}$ gives continuation in date-$n$ currency,
matching $C_n$ in the earlier recursion. These targets are observed simulated payments;
the network learns their conditional average. Its exercise decision compares
$G_n$ with $c_\theta(S_n)$. Working backward makes the later rule available
before fitting the current one. This is the continuation-regression approach
used in the 2020 paper~\cite{american}.

The 2019 direct approach instead learns an exercise decision. During training,
a sigmoid output $f_\theta(s)$ lies between zero and one. At each fitted date,
it maximizes the average of
\[
 G_n f_\theta(S_n)+(1-f_\theta(S_n))G_{\tau_{n+1}}.
\]
The objective rewards a larger immediate-payment weight in states where
exercise is preferable, and a larger continuation weight where waiting is
preferable. A threshold converts the selected function into a binary exercise
rule. Both algorithms use future simulated outcomes to improve a rule that,
when applied, observes only the current state~\cite{stopping}.

A state is \emph{Markov} when its current value determines the conditional
law of the next state, given the model. Stock price supplies that state in GBM;
Heston uses stock and variance together. When variance or a market regime is
only partly observed, estimates and relevant history can enter the policy
observation. A recurrent network can summarize that history through an
internal state. The particular reset-LSTM engine in the completed stock-and-option
study instead shares weights across dates and uses financial holdings and
wealth to carry its specified history.

The writer's network has a different output: target stock or instrument
holdings. Write $h_n=\pi_\theta(x_n,h_{n-1},w_n)$, where $x_n$ contains the
market observations, $h_{n-1}$ contains previous holdings and $w_n$ is wealth
when that input is enabled. The account calculation converts all decisions
into a final profit $\Pi_T(\theta)$. Training can minimize
\[
 J(\theta)=\E[\Pi_T(\theta)^2].
\]
The holder's actual settlement, financing and transaction fees determine that
profit. The episode begins with the received premium and zero hedge inventory.
At each hedge date the action is the new vector of continuous holdings. The
market simulator advances the observations, and the account advances cash and
positions. The episode settles at the holder's exercise date, then carries
residual cash to the reporting horizon. For squared-loss training its terminal
reward can be written $-\Pi_T^2$. A CVaR objective depends on the loss
distribution across episodes, which the batch calculation supplies. An analytic delta hedge can be a comparison strategy; the writer's
objective obtains its feedback from the simulated account.

Different risk preferences change the objective. Mean squared error penalizes
both positive and negative departures from the chosen capital target. A
downside squared loss uses $\E[((-\Pi_T)^+)^2]$. For loss $L=-\Pi_T$, a
CVaR objective at confidence level $\alpha$ can be written
\[
 \min_{\eta}\left\{\eta+
       \frac{1}{1-\alpha}\E[(L-\eta)^+]\right\}.
\]
CVaR averages the worst tail of losses; $\eta$ is a fitted threshold. An
entropic objective $(1/\gamma)\log\E[e^{\gamma L}]$, with risk-aversion
parameter $\gamma>0$, assigns increasing weight to large losses. The selected
objective and premium determine what the learner is trying to improve.
For example, $\E[\Pi_T^2]=\operatorname{Var}(\Pi_T)+(\E[\Pi_T])^2$.
Mean squared profit therefore measures dispersion and the squared mean departure
from the chosen zero-residual target. Its units are currency squared; positive
average profit is a separate reported statistic.

This is a sequential control problem: state, action, transition and eventual
outcome. Reinforcement learning studies how to choose actions from such
experience~\cite{sutton}. Our writer uses a differentiable simulator and
optimizes complete paths, propagating derivatives through the trading account.
The continuation learner uses regression targets, and the direct holder
optimizes a stopping reward. Other Jormungandr policy engines can implement
actor--critic or value-based RL algorithms. The method name in a study should
identify which of these calculations produced its decisions.

A trained policy is a function of its supplied observations. A new stock
price calls for a new evaluation of that function. The policy also carries
training assumptions about liability, market law, horizon, costs and units.
A policy intended to serve arbitrary portfolios needs those portfolio
characteristics represented in its inputs and training distribution. Two
portfolios with identical pathwise payments, initial capital, market law, risk
objective and admissible hedging/execution terms share the same economic objective; matching their current Greeks alone
provides a local comparison of exposure.

\section{Automatic differentiation}
Training needs to know how changing a weight changes the objective. A derivative
answers that local question. \emph{Automatic differentiation}, or AD, applies
the chain rule to the actual sequence of numerical operations. Reverse mode
is also called adjoint algorithmic differentiation, or AAD; neural-network
backpropagation uses the same principle~\cite{ad}.

Consider the small calculation $a=wx$, $b=\tanh(a)$ and
$\ell=(b-y)^2$. Here $x$ is an observed input, $w$ is a weight and $y$ is a
target. The chain rule gives
\[
 \frac{\partial\ell}{\partial w}
 =2(b-y)(1-b^2)x.
\]
For a large network, the program applies the same local rules through many
operations. An optimizer such as Adam uses the resulting gradients, together
with running first and second moments, to update the weights. Learning rate,
normalization and the units of inputs influence the numerical trajectory.

\begin{center}\begin{minipage}{\linewidth}\centering
\begin{tikzpicture}[font=\normalsize,>=Latex,x=1cm,y=1cm]
\node (w) at(0,0){$w,x$};\node (a) at(1.65,0){$a=wx$};
\node (b) at(3.7,0){$b=\tanh a$};\node (l) at(6.05,0){$\ell=(b-y)^2$};
\draw[->] (w)--(a);\draw[->] (a)--(b);\draw[->] (b)--(l);
\draw[->,darkblue] (5.8,-.6)--(.5,-.6);
\node[darkblue,anchor=north] at(3.15,-.65){Reverse pass: accumulate chain-rule derivatives};
\end{tikzpicture}
\captionof{figure}{AD follows the numerical computation. In hedge training,
the intermediate operations include trades, costs, financing and settlement.}
\end{minipage}\end{center}

A pricing sensitivity and a training gradient differ in what is varied.
For delta, the input being differentiated is the stock price; model parameters
and the contract's valuation convention are held as specified. For a writer
update, the differentiated inputs are policy weights. The same AD tooling
can serve both tasks, provided the intended inputs remain connected to the
calculation graph.

In forward mode, derivatives travel with the evaluated quantities. This is
useful for a small number of input directions or selected higher-order
quantities. Reverse mode propagates derivatives back from a scalar output
to many inputs and usually retains intermediate values, or recomputes them,
for that backward calculation. Those saved values contribute to memory use.
A directional finite difference, obtained by perturbing an input and
re-evaluating, supplies an independent numerical check on a derivative.

Financial functions include kinks. The put payoff changes slope at its strike;
a proportional cost uses an absolute value; an exercise decision changes a
branch. The estimator must respect that structure. Differentiating a sample
path payoff can estimate certain sensitivities under appropriate regularity,
while a second derivative of the same sampled kink needs additional treatment.
For example, write a GBM terminal stock as $S_T=S_0A$ with random multiplier
$A$. The put payment $(K-S_0A)^+$ is piecewise linear in $S_0$, so its ordinary
second derivative is zero away from the kink on each path. The expected
option value can nevertheless have positive gamma. A suitable sensitivity
estimator accounts for how probability mass near the strike changes with $S_0$.
Volt checks its supported gamma estimators separately. The learned
holder is held fixed during the writer's stopped-account training, and the
account removes later policy actions after settlement.

A reliable training calculation also checks its scale. Stock holdings may be
fractions, while another hedge instrument may require thousands of quoted
units. Explicit input transforms and action scales can improve optimization
while preserving the converted economic holdings. A good comparison restores
all results to their declared financial units before comparing profits,
fees and exposures. Our Heston studies show that these implementation choices
can materially affect the fitted policy.

\section{Validation and completed experiments}\label{sec:evidence}
A policy's training loss says how it performed on the paths used to improve
it. The practical question is how its selected version behaves on different
paths from the intended application law. We assess profit distributions,
loss objectives, trading costs and comparison hedges on common fresh paths.
An independent cash calculation checks the financial arithmetic, and separate
derivative tests check that optimization receives the intended gradient.

Valuation of a learned holder offers another useful check. Under $\Q$, its
expected discounted exercise payment is a lower value relative to the best
admissible exercise rule. A Monte Carlo average estimates that expectation
with sampling uncertainty. An upper-value construction uses a martingale
$M_n$, starting at zero, to form
\[
 U=\E^{\Q}\!\left[\max_{n\text{ allowed}}(G_n-M_n)\right].
\]
The martingale has zero expected gain when sampled at any allowed stopping
time under the stated finite-horizon assumptions. Taking the maximum over
allowed dates then makes this expectation at least the optimal stopping value.

Our initial estimator builds increments from a selected time-zero-discounted
value approximation
$v_{n+1}$ and independent conditional simulations:
\[
 \delta M_{n+1}=v_{n+1}(S_{n+1})-
     \frac1K\sum_{k=1}^{K}v_{n+1}(\widetilde S_{n+1}^{(k)}).
\]
Both terms have the same conditional mean given the current state. The fresh
conditional children and the next outer stock draw preserve the martingale
property with the stated information ordering. The papers describe a
policy-tail construction; our continuation-proxy estimator supplies a
separately identified alternative. More inner draws can reduce its extra
simulation noise. The upper sample average has uncertainty too, so the
report includes intervals around both expectations.

Volt's numerical pricing engine supports American exercise through a
finite-difference obstacle calculation, alongside European and dated Bermudan
rights. Its tests include tree and QuantLib comparisons, dated cash dividends
and mesh refinements. The learning results below use specified Bermudan dates.
Their holder decisions, prices and hedge errors are evaluated under those
calendars; refining the exercise dates provides a route toward continuous
exercise. The pricing and learning guides document the separate
interfaces~\cite{voltapi}.

Our one-stock application uses a put with spot and strike 100, 20\% GBM
volatility and 5\% interest over one year. Two experiment calendars answer
different questions. In the writer study, exercise is quarterly and hedging
is monthly. Three seeds per cost setting each train for 3,000 updates using
16,384 training and 4,096 validation paths. The six selected policies are
then compared on 16,384 fresh common paths.

\begin{center}\begin{minipage}{\linewidth}\centering
\begin{tabular}{lrr}
\toprule
Strategy & Zero stock fees & 10 bp stock fees\\
\midrule
Learned writer & 4.013 & 4.079\\
Numerical delta & 3.935 & 3.984\\
Selected stock band & 3.937 & 3.986\\
Unhedged & 61.443 & 61.443\\
\bottomrule
\end{tabular}
\captionof{table}{Mean squared terminal-bank loss in currency squared for one underlying unit.
One basis point is 0.01\%; 10 bp is a fee fraction of 0.001 of traded notional.
Learned entries average three selected policies.}
\end{minipage}\end{center}

The learned writer removes about 93\% of unhedged squared loss. Numerical
delta retains an advantage of about 2\% of its own squared loss at this budget.
The matching comparison intervals and all cash-account audits are given in
the application paper. These values describe the stated synthetic model,
contract, capital and execution assumptions.

The holder study permits exercise on all twelve monthly dates. It trains
three seeds for each learned-stopping method, then evaluates discounted
payments on 65,536 fresh paths. Payment estimates range from 5.952 to 5.983;
the alternate upper estimates range from 6.262 to 6.292 using 4,096 outer
paths and 64 conditional children per step. All six asymptotic 95\% price
enclosures contain the matching finite-difference value of 6.0428. The
approximately 0.3 estimator gap measures remaining numerical room in this
particular comparison.

The separate 2020 interval-hedge calculation learns how to hedge changes in
approximate continuation values from one exercise date to the next. Three
selected hedges have discounted-error MSEs of 3.986--4.116 on 32,768 additional
paths, versus 54.57--58.60 for cash alone at their declared initial capital.
The full writer and interval methods use different targets, financing
representations and assessment cohorts, so their reports identify those
conventions. The papers' published multidimensional max-call experiments
provide the next replication benchmark.

\section{Portfolios and factor models}
A portfolio adds many payments and overlapping risks. For a terminal-only
book, its liability can be written $Z=\sum_j n_jm_jg_j$, where $n_j$ is signed
liability quantity, positive for a payment owed by the writer, and $m_j$ is
its multiplier. Multiple exercise dates require
that same aggregation at the actual settlement dates. The hedge policy can
choose a vector of instrument holdings, and the cash account adds the gains
and costs of each instrument.

A \emph{factor model} describes shared movements across assets. One useful
return representation is
\[
 r_{n+1}=\mu_n+B_n f_{n+1}+\varepsilon_{n+1}.
\]
Here $r_{n+1}$ is a vector of returns, $\mu_n$ is their conditional mean,
$B_n$ is a loading matrix, $f_{n+1}$ is a factor shock and $\varepsilon_{n+1}$
is an idiosyncratic residual. Factor shocks and residuals have conditional
mean zero here; any predictable means enter $\mu_n$. A specified shock law and residual law turn the
representation into a return generator. With log returns, prices evolve as
$S_{n+1}^i=S_n^i\exp(r_{n+1}^i)$, preserving positive prices. Simple returns
instead require $S_{n+1}^i=S_n^i(1+r_{n+1}^i)$ and admissible return support.

Principal component analysis, or PCA, finds orthogonal directions that explain
sample variation~\cite{pca}. For a return covariance estimate $\Sigma$, its eigenvectors
supply these directions and its eigenvalues give their variances. Retaining
a few components compresses shared covariance. Generative use additionally
specifies how shocks evolve, how residuals behave and how estimates update.
Forecasting the conditional mean requires a further predictive model and
an assessment of its uncertainty.

There are several research choices before attaching factors to a hedge.
An economically named sector exposure, a PCA return component and a forecast
signal represent different information. Their definitions, observation times,
estimation windows and response to regime changes need to be specified.
Option observations add implied volatility, skew and term structure, with
liquidity and quote quality influencing the cross-section. A factor formed
from Greeks also depends on each contract's pricing conventions and units.

These quantities can also organize a live risk view. If $b_{ik}$ is stock
$i$'s loading on a small dimensionless factor return shock $f_k$, its approximate
stock-price change is $S_i b_{ik}f_k$. For signed option units $u_j$, positive for long positions, a first-order
portfolio exposure to that factor is
\[
 A_k=\sum_j u_j\Delta_j S_{i(j)}b_{i(j)k},
\]
where $i(j)$ identifies contract $j$'s underlying. The linear price response is
approximately $A_k f_k$. Volatility factors add vega times their own volatility
loadings, with consistent decimal or percentage-point units. Gamma contributes
curvature. This gives stock, sector and factor views of the same option book,
provided its positions and model conventions are known.

The intended experiment keeps the liability, paths, costs and assessment
cohort aligned while comparing a hedge-only policy, a signal-trading policy
and a combined policy. This reveals which gains come from protecting the
liability and which depend on predicted returns. The signal law belongs to
$\PP$; coherent option valuation still needs its pricing specification.
Research on removing predictable drift from hedge training gives a related
way to study this distinction~\cite{drift}.

Factor research therefore precedes the empirical combination. We will examine
what existing models estimate, how to generate admissible joint paths, and
which forecasts remain useful on data observed strictly after fitting. The
published max-call benchmark first gives us a common multi-asset pricing and
exercise reference under a fully declared model.

\section{Applications and computational execution}
The financial equations become an application through a few connected tools.
Volt holds contract and portfolio definitions, generates or accepts scenarios,
calculates prices and Greeks, and evaluates the hedge account. Burn supplies
tensor operations, automatic differentiation and optimization. CubeCL supports
GPU execution beneath the selected backend. Jormungandr supervises a persistent
learner and supplies its HTTP lifecycle, dataset and inference operations.
Hypercube publishes dated market and risk records for analytical applications.
The implementation uses the backends and precision supported by its pinned
build; the upstream Burn documentation explains the underlying
architecture~\cite{burn}.

\begin{center}\begin{minipage}{\linewidth}\centering
\begin{tikzpicture}[font=\normalsize,>=Latex,
 block/.style={draw=darkblue,rounded corners,align=center,inner sep=4pt}]
\node[block,text width=2.5cm] (api) at(0,0){Jormungandr\\HTTP and model lifecycle};
\node[block,text width=2.8cm] (volt) at(3.55,0){Volt worker\\Policy and\\financed account};
\node[block,text width=2.8cm] (burn) at(3.55,-1.4){Burn tensors / AD\\CPU or\\CubeCL / GPU};
\node[block,text width=2.5cm] (cube) at(0,-2.7){Hypercube\\Dated observations\\and risk records};
\draw[->] (api)--(volt);\draw[->] (volt)--(burn);
\draw[->] (cube.north)--node[left,align=right]{Application\\observations}(api.south);
\draw[->,dashed] (volt.south west)--node[below,sloped]{Calculated risk}(cube.east);
\end{tikzpicture}
\captionof{figure}{Application roles. The market-data path and the resident
training calculation can operate on different schedules.}
\end{minipage}\end{center}

A GPU performs many related arithmetic operations in parallel. A batch of
paths becomes matrices and vectors: a tensor is simply an array with a
specified shape and data type. For example, stock marks have dimensions
$[\text{path},\text{date},\text{instrument}]$. Keeping tensors resident on
the device lets several calculations reuse them before transferring results
back to the application. A Rust array uploaded through an API may involve
allocation and copying even on a machine with physically shared memory;
the selected tensor interface defines ownership and synchronization.

\emph{Memory budgeting} means estimating the peak storage needed before
admitting a workload. With $B$ paths, $L$ dates, $d$ instruments and $s$ bytes
per scalar, marks alone use approximately $BLds$ bytes. At $B=100{,}000$,
$L=31$, $d=10$ and f32 storage, that is 124 million bytes, about 118 MiB.
Features, activations, gradients, optimizer moments, temporary kernels and
transactional copies add to this amount. Streaming smaller batches keeps
the study's total path count independent of its resident batch size.

GPUs help when enough parallel work amortizes dispatch, synchronization and
transfer time. A small date-specific network can run faster on a CPU. Fusion
can combine compatible tensor operations into fewer kernels and reduce
intermediate traffic. Floating-point precision and reduction order also
affect numerical results, especially when subtracting nearby values. We
measure the actual configured training or pricing calculation, with its
accuracy checks, to choose a device.

Three CPU commands demonstrate the learning application:
\end{SNcolumns}
\noindent\begin{minipage}{\textwidth}
\begin{lstlisting}[basicstyle=\ttfamily\normalsize,frame=single,breaklines=true]
cargo build --release -p volt-compute --bin deep_stopping --bin bermudan_market --bin native_learner
python3 examples/deep_stopping.py --output /tmp/volt-stopping-example
python3 examples/bermudan_learning.py --output /tmp/volt-bermudan-example
\end{lstlisting}
\end{minipage}
\begin{SNcolumns}
The first example fits a holder and an interval hedge, then estimates valuation
and hedging results on independent paths. The second fits a writer against a
numerical holder, selects its checkpoint on validation and compares it with
a financed delta control on fresh paths. Each retains the inputs and results
needed to inspect the complete sequence. The default examples use small,
bounded budgets; the completed studies above have their own fixed protocols.

Through Jormungandr, an application creates a native model, registers dataset
partitions, uploads a scenario batch, requests a versioned update, evaluates
held-out paths and retrieves a checkpoint or target holdings. The writer's
batch supplies stock marks, decision-time features and actual settlement
indices and amounts. Optional income and initial-capital arrays make the
financial convention explicit. Current observations and holdings then feed
inference until the contract settles. The API guides provide exact request
schemas and executable examples~\cite{voltapi,jormapi}.

\section{Exercises and further reading}
The following short checks connect the main ideas. They are useful before
changing a market model or interpreting a training result.

\begin{enumerate}[leftmargin=1.5em]
\item\begin{minipage}[t]{\linewidth}A put with strike 100 expires at stock 92. For three contracts with
multiplier 100, what is the payment? It is $3\times100\times8=2{,}400$ to
the holder. Premium and hedge gains determine the writer's final profit.
\end{minipage}\par
\item\begin{minipage}[t]{\linewidth}A model vega is 24 per underlying unit for volatility expressed as a
decimal. What is the approximate effect of a one-percentage-point increase
on one long contract with multiplier 100? It is $24\times0.01\times100=24$.
\end{minipage}\par
\item\begin{minipage}[t]{\linewidth}A portfolio's Monte Carlo standard error is 0.04 using 10,000 independent
paths. How many comparable paths target 0.02? Approximately 40,000, keeping
the same distribution and estimator.
\end{minipage}\par
\item\begin{minipage}[t]{\linewidth}On an exercise date, a put's immediate payment is 7 and its model
continuation value is 6.5. What does the value-maximizing holder do? It
exercises. The writer closes its hedge and settles according to the event order.
\end{minipage}\par
\item\begin{minipage}[t]{\linewidth}A factor model explains most historical covariance. What must be added
to generate return scenarios? A shock distribution and evolution law,
a residual model, the conditional-mean convention and a dated estimation rule.
\end{minipage}\par
\item\begin{minipage}[t]{\linewidth}A policy responds correctly to a new stock price. What determines whether
it can handle a different portfolio? Its observations and training distribution
must represent the relevant portfolio payments, quantities and contract terms.
\end{minipage}\par
\end{enumerate}

For contract terms and everyday quotation conventions, begin with the Options
Industry Council material~\cite{oicpricing,oicexercise}. Then read the worked
accounts in this primer and the Jormungandr application paper. Glasserman's
\emph{Monte Carlo Methods in Financial Engineering} develops simulation,
estimation and variance reduction; Shreve's \emph{Stochastic Calculus for
Finance II} develops pricing measures and continuous-time models~\cite{glasserman,shreve}.

Read the original Deep Hedging paper next for after-cost portfolio objectives
and learned trading decisions~\cite{deephedging}. Becker, Cheridito and Jentzen's
2019 paper develops direct stopping, followed by their 2020 treatment of
continuation regression, bounds and hedging~\cite{stopping,american}. Baydin
and colleagues explain AD and its relationship to backpropagation~\cite{ad}.
Sutton and Barto provide the wider sequential-learning framework~\cite{sutton}.
Jolliffe and Cadima review PCA and its interpretation before the factor-model
extensions~\cite{pca}. The repository guides then connect these methods to the
supplied APIs.

A useful progression is to reproduce the one-period cash arithmetic, run both
small examples, inspect a saved exercise path and compare a selected hedge
on fresh scenarios. From there, study the common multi-asset benchmark before
introducing estimated factors and predictive signals. Each step adds a
specific source of complexity while retaining calculations that can be checked.
\end{SNcolumns}
\clearpage
\begin{SNcolumns}
\begin{thebibliography}{99}
\bibitem{oicpricing} Options Industry Council. Options Pricing. Contract values and pricing inputs. \url{https://www.optionseducation.org/optionsoverview/options-pricing}.
\bibitem{oicexercise} Options Industry Council. Exercising Options. \url{https://www.optionseducation.org/optionsoverview/exercising-options}.
\bibitem{glasserman} P. Glasserman. \emph{Monte Carlo Methods in Financial Engineering}. Springer, 2003. \url{https://doi.org/10.1007/978-0-387-21617-1}.
\bibitem{shreve} S. E. Shreve. \emph{Stochastic Calculus for Finance II: Continuous-Time Models}. Springer, 2004. \url{https://link.springer.com/book/9780387401010}.
\bibitem{heston} S. L. Heston. A Closed-Form Solution for Options with Stochastic Volatility with Applications to Bond and Currency Options. \emph{Review of Financial Studies} 6(2):327--343, 1993. \url{https://doi.org/10.1093/rfs/6.2.327}.
\bibitem{deephedging} H. Buehler, L. Gonon, J. Teichmann and B. Wood. Deep Hedging. 2018. \url{https://arxiv.org/abs/1802.03042}.
\bibitem{stopping} S. Becker, P. Cheridito and A. Jentzen. Deep Optimal Stopping. \emph{JMLR} 20(74):1--25, 2019. \url{https://jmlr.org/papers/v20/18-232.html}.
\bibitem{american} S. Becker, P. Cheridito and A. Jentzen. Pricing and hedging American-style options with deep learning. \emph{JRFM} 13(7):158, 2020. \url{https://arxiv.org/abs/1912.11060}.
\bibitem{ad} A. G. Baydin, B. A. Pearlmutter, A. A. Radul and J. M. Siskind. Automatic differentiation in machine learning: a survey. \emph{JMLR} 18(153):1--43, 2018. \url{https://arxiv.org/abs/1502.05767}.
\bibitem{sutton} R. S. Sutton and A. G. Barto. \emph{Reinforcement Learning: An Introduction}, second edition. MIT Press, 2018. \url{https://mitpress.mit.edu/9780262039246/reinforcement-learning/}.
\bibitem{drift} H. Buehler, P. Murray, M. S. Pakkanen and B. Wood. Deep Hedging: Learning to Remove the Drift under Trading Frictions with Minimal Equivalent Near-Martingale Measures. 2021. \url{https://arxiv.org/abs/2111.07844}.
\bibitem{pca} I. T. Jolliffe and J. Cadima. Principal component analysis: a review and recent developments. \emph{Philosophical Transactions A} 374:20150202, 2016. \url{https://pmc.ncbi.nlm.nih.gov/articles/PMC4792409/}.
\bibitem{burn} Burn project. Backend and automatic-differentiation documentation. \url{https://burn.dev/books/burn/}.
\bibitem{voltapi} Volt. Exercise pricing, native hedge learner, Bermudan scenarios and learned stopping guides. \url{https://github.com/ingvor-restrat/volt/blob/main/docs/exercise-pricing.md}. \url{https://github.com/ingvor-restrat/volt/blob/main/crates/volt-compute/docs-native-learner.md}. \url{https://github.com/ingvor-restrat/volt/blob/main/docs/bermudan-market.md}. \url{https://github.com/ingvor-restrat/volt/blob/main/crates/volt-compute/docs-deep-stopping.md}.
\bibitem{jormapi} Jormungandr. Native learner API and mathematical application tutorial. \url{https://github.com/ingvor-restrat/jormungandr/blob/main/docs/native-learners.md}. \url{https://github.com/ingvor-restrat/jormungandr/tree/main/docs/latex/hedging}.
\end{thebibliography}
\end{SNcolumns}
\end{document}
